\section*{Chapter \ref{ch:g6:health-and-population} --- Longer Lives, Smaller Families}

\begin{solution}{exo:g6:health-and-population:1}
About 34 years in 1900 and about 73 years in 2024 (the source gives $33.6$ and
$73.5$).
\end{solution}

\begin{solution}{exo:g6:health-and-population:2}
The fewest: East Asia and Pacific ($1.34$ \omterm{def:g6:health-and-population:births-per-woman}{births per woman}). The most:
Sub-Saharan Africa ($4.26$).
\end{solution}

\begin{solution}{exo:g6:health-and-population:3}
(a) $\frac{37}{1\,000} = 3.7\,\%$. (b) $50\,000 \times \frac{37}{1\,000} = 1\,850$
children.
\end{solution}

\begin{solution}{exo:g6:health-and-population:4}
(a) Most dots near 2\,000 dollars lie between about 40 and 100 deaths per
1\,000 births. (b) Near 50\,000 dollars, the dots lie below about 5.
\end{solution}

\begin{solution}{exo:g6:health-and-population:5}
(a) $0 + 0 + 0 + 2 + 4 + 40 + 55 + 60 + 68 + 71 = 300$, and $300 \div 10 = 30$
years. (b) None: five died before five, five after forty. (c) $(40 + 55 + 60 + 68
+ 71) \div 5 = 294 \div 5 = 58.8$ years. A life expectancy of 30 years does not
mean that people died at 30: it is pulled down by many deaths of young children,
while those who passed childhood often lived to 50, 60 or more.
\end{solution}

\begin{solution}{exo:g6:health-and-population:6}
$1\,200 \div 4\,800 \times 100 = 25$: 25 people aged 65 and over for every 100
of working age.
\end{solution}

\begin{solution}{exo:g6:health-and-population:7}
In the early 1990s (the table gives 1993).
\end{solution}

\begin{solution}{exo:g6:health-and-population:8}
(a) Two dots where the one further right (richer) is also higher (more child
deaths). (b) Yes, many: for example Equatorial Guinea, richer than Sri Lanka, has a
\omterm{def:g6:health-and-population:child-mortality}{child mortality rate} more than ten times as high. (c) ``Richer countries
\emph{usually} have lower child mortality: on average, the rate falls as \omterm{def:g6:living-standards:income-per-person}{income
per person} rises, but countries with the same income can be far apart.''
\end{solution}

\begin{solution}{exo:g6:health-and-population:9}
(a) $5\,000 + 150 - 60 = 5\,090$; $5\,090 + 140 - 65 = 5\,165$; $5\,165 + 130 - 70 =
5\,225$. (b) It is still growing, but more slowly each year: by 90, then 75,
then 60 people.
\end{solution}

\begin{solution}{exo:g6:health-and-population:10}
(a) $24\,230 - 24\,203 = 27$ dollars. (b) $30.6 \div 2.8 \approx 10.9$ times.
(c) For example: access to clean water and sanitation; vaccination; the number of
clinics and midwives; how long mothers went to school; how income is shared.
\end{solution}

\begin{solution}{exo:g6:health-and-population:11}
Now: $2\,000 \div 6\,000 \times 100 \approx 33.3$. After 5 years: $2\,300 \div 5\,800
\times 100 \approx 39.7$. After 10: $2\,600 \div 5\,600 \times 100 \approx 46.4$.
After 15: $2\,900 \div 5\,400 \times 100 \approx 53.7$. After 20: $3\,200 \div 5\,200
\times 100 \approx 61.5$. The rises are $6.4$, $6.7$, $7.3$ and $7.8$: not the same;
the ratio rises faster and faster, because the top grows while the bottom shrinks.
\end{solution}

\begin{solution}{exo:g6:health-and-population:12}
(a) For example: men and women; people in towns and in the countryside; people
with more and with less schooling; the richest and the poorest tenth. (b) Death
counts and population counts by age for each group separately, to build a life
table for each. (c) No: the chart has one number per region; it says nothing
about groups inside a country.
\end{solution}

\begin{solution}{pb:g6:health-and-population:1}
\textbf{1.} 10 children.
\textbf{2.} $10 \div 40 \times 1\,000 = 250$ per 1\,000 births: one child in four.
(It is an estimate: some of the children who died were born in earlier years.)
\textbf{3.} The children's ages add up to $0 \times 5 + 1 + 1 + 2 + 3 + 4 = 11$, the
others' to $15 + 28 + 35 + 42 + 50 + 58 + 63 + 67 + 72 + 81 = 511$. The average is
$(11 + 511) \div 20 = 522 \div 20 = 26.1$ years.
\textbf{4.} $511 \div 10 = 51.1$ years. The average of question 3 is pulled down by
the ten children who died very young; those who survived childhood lived about
twice as long.
\textbf{5.} $800 + 10 \times (40 - 20) = 1\,000$ people.
\textbf{6.} $1\,000 + 10 \times (40 - 12) = 1\,280$.
\textbf{7.} $1\,280 + 10 \times (30 - 12) = 1\,460$.
\textbf{8.} $(1\,460 - 800) \div 800 = 660 \div 800 = 0.825$: a rise of $82.5\,\%$.
It grew fastest from 1835 to 1845 (280 people): deaths had fallen while births
had not yet fallen.
\textbf{9.} $380 \div 900 \times 100 \approx 42.2$.
\textbf{10.} $480 \div 800 \times 100 = 60.0$.
\textbf{11.} Today $900 \div 380 \approx 2.37$ people of working age for each person
aged 65 and over; in twenty years $800 \div 480 \approx 1.67$.
\textbf{12.} For example: fewer working people will pay the taxes that fund care
homes, clinics and pensions for more older people; the town may need more care
workers and fewer school places.
\end{solution}
